\chapter{双边遥操作的绝对稳定性与时延（Llewellyn 判据）}\label{ch:teleop-stability}

% ════════════════════════════════════════════════════════════════
%  机械臂部·遥操作子部 第 2 章：承 ch:teleop-twoport（h 参数/H_ideal），
%  引出 ch:teleop-wave（波变量）。理论吸收自 Llewellyn 1952 / Lawrence 1993 /
%  Hannaford 1989 / Colgate-Schenkel 1997，源稿 Tutorial D05。
%  记号归一：能量端口 P=f^T v（与 \cref{eq:os-port-power} 一致），
%  机械臂动力学 M q̈+C q̇+g=tau，静力对偶 tau=Jᵀf。
% ════════════════════════════════════════════════════════════════

\begin{practice}[前置自测]
开始之前，先试答下列问题。若能流畅作答，可只速读本章的判据陈述与速查表；若卡住，正是本章要为你打通的。
\begin{enumerate}
  \item 一个单端口机械系统“被动（passive）”的能量定义是什么？写出端口功率积分的不等式。
  \item 二端口混合参数 $\mathbf H$ 的四个元 $h_{11},h_{12},h_{21},h_{22}$ 各自的因果含义是什么？理想透明度对应的 $\mathbf H_{\rm ideal}$ 长什么样？
  \item “两个各自被动的子系统，负反馈互联后一定仍然稳定吗？”——若不一定，缺了什么条件？
  \item 通信往返时延 $T$ 会给力/位置回路的开环传递函数乘上一个什么因子？它对相位做了什么？
  \item 数字实现一面“虚拟墙”时，采样周期 $T$ 越大，可稳定渲染的刚度上限是变高还是变低？为什么？
\end{enumerate}
\end{practice}

\section*{本章目标}
学完本章，你应当能够：
\begin{enumerate}
  \item \textbf{说清}为何双边遥操作必须追求\emph{绝对}（而非名义）稳定——操作者与环境阻抗会在极端区间任意游走；
  \item \textbf{陈述并解读} Llewellyn 四判据 $\mathrm{L1}$--$\mathrm{L4}$，把每条对应到“端口被动”与“互耦合不创生能量”的物理含义；
  \item \textbf{完整推导}闭环特征方程 $(Z_h+h_{11})(1+h_{22}Z_e)=h_{12}h_{21}Z_e$，并由极端终端 + 一般被动终端导出 $\mathrm{L1}$--$\mathrm{L4}$ 与稳定因子 $\eta(\omega)$；
  \item \textbf{独立计算}一个 1-DOF PD 位置跟踪 + 力反射架构的 $h$ 参数与 $\eta(\omega)$，并判稳；
  \item \textbf{证明}透明度-稳定性的根本矛盾：理想透明 $\Rightarrow\eta\equiv0$（零鲁棒裕度），加本地阻尼 $\Rightarrow\eta'=2b_mb_s>0$ 但 $Z_{to}\neq Z_e$；
  \item \textbf{定量分析}时延：$h_{12}h_{21}=-|C_fC_p|e^{-j\omega T}$ 如何在 $\omega=\pi/T$ 把 $\eta$ 拉负、阻尼要求升为 $b_mb_s\ge|C_fC_p|$、为何 $>200\,$ms 时 Llewellyn 几近无解；
  \item \textbf{复用}本书 Colgate--Schenkel 采样无源界（\cref{eq:os-discrete-passive}）解释 $Z$-width 与“ALOHA $50\,$Hz 力反馈不可行”。
\end{enumerate}

\subsection*{本章知识导航}
本章是一条“\emph{把‘想让操作者真实感受远端’的渴望，约束在‘对任意人/环境都不失稳’的红线之内}”的主线。结构如下：

\begin{center}
\small
\begin{tikzpicture}[
  node distance=4.5mm and 7mm, >=Stealth,
  blk/.style={draw, rounded corners, align=center, font=\small, inner sep=3pt, fill=main!6, draw=main!60!black},
  grp/.style={draw, rounded corners, align=center, font=\small, inner sep=3pt, fill=second!6, draw=second!60!black},
  ln/.style={->, thick, gray!60!black}]
\node[blk] (motiv) {为何要\\\emph{绝对}稳定};
\node[blk, right=of motiv] (crit) {Llewellyn\\$\mathrm{L1}$--$\mathrm{L4}$};
\node[blk, right=of crit] (proof) {闭环特征方程\\$\to\eta(\omega)$};
\node[grp, below=8mm of motiv] (ex) {1-DOF\\worked example};
\node[grp, right=of ex] (trade) {透明度--稳定\\trade-off ($\eta{=}0$)};
\node[grp, right=of trade] (delay) {时延\\$e^{-j\omega T}$};
\node[blk, below=8mm of trade, fill=third!8, draw=third!60!black] (zwidth) {$Z$-width\\Colgate--Schenkel};
\node[blk, right=of zwidth, fill=third!8, draw=third!60!black] (wave) {波变量\\（下一章）};
\draw[ln] (motiv)--(crit); \draw[ln] (crit)--(proof);
\draw[ln] (proof.south) -- ++(0,-4.5mm) -| (ex.north);
\draw[ln] (ex)--(trade); \draw[ln] (trade)--(delay);
\draw[ln] (delay.south) |- (zwidth.east);
\draw[ln] (zwidth)--(wave);
\end{tikzpicture}
\end{center}

\noindent\textbf{推荐路径}：第 \ref{sec:ts-why}--\ref{sec:ts-proof} 节是主干（动机 $\to$ 四判据 $\to$ 证明），任何要做力反馈系统的读者必读；第 \ref{sec:ts-example} 节把判据落成一次手算；第 \ref{sec:ts-tradeoff}--\ref{sec:ts-delay} 节是本章的“坏消息”（透明与稳定不可兼得、时延雪上加霜），它们正是下一章波变量（\cref{ch:teleop-wave}）的动机；第 \ref{sec:ts-zwidth} 节把离散采样的影响接回本书已证的 Colgate--Schenkel 界。

\subsection*{前置知识桥接}
本章紧承上一章二端口网络（\cref{ch:teleop-twoport}）：那里建立了端口变量（力 $f$／速度 $v$）、混合参数矩阵 $\mathbf H$、操作者感受阻抗 $Z_{to}$，以及理想透明度 $\mathbf H_{\rm ideal}=\big[\begin{smallmatrix}0&1\\-1&0\end{smallmatrix}\big]$。本章要回答上一章\emph{回避}的问题：这套 $h$ 参数在什么条件下、对\emph{任意}操作者与环境都不会失稳。无源性、正实/KYP 与采样数据无源的\emph{通用理论}已在操作空间一章（\cref{ch:operational-space}）证过（\cref{thm:os-kyp,eq:os-discrete-passive}），本章只取遥操作所需、不重证。能量端口约定 $P=\mathbf f^\top\mathbf v$ 与 \cref{eq:os-port-power} 一致。

\subsection*{如果跳过本章会怎样}
\begin{itemize}
  \item 你会调出一个在实验室“看起来稳”的力反馈遥操作系统，却无法回答“换一个更紧的握法、撞上更硬的墙，它还稳吗”——而手术机器人等场景里，这个问题的答案关乎安全；
  \item 你会把“透明度上不去”归咎于“控制器不够好”，反复加增益直到振荡，而看不见 $\eta=0$ 这条由物理定律划定的天花板；
  \item 你会不理解为什么 ALOHA 这类 $50\,$Hz 的遥操作明明便宜地装个力传感器就行，却\emph{故意}不做力反馈——\cref{sec:ts-zwidth} 给出定量理由。
\end{itemize}

\begin{note}
\textbf{本章记号与约定.}\;
沿用 \cref{ch:teleop-twoport} 的混合参数：端口 1（master，操作者侧）取阻抗因果，端口 2（slave，环境侧）取导纳因果，
\begin{equation}
  \begin{bmatrix}F_h\\-V_s\end{bmatrix}=\mathbf H\begin{bmatrix}V_m\\F_e\end{bmatrix},\qquad
  \mathbf H=\begin{bmatrix}h_{11}&h_{12}\\h_{21}&h_{22}\end{bmatrix}.
  \label{eq:ts-hparam}
\end{equation}
其中 $V_s$ 带负号是\emph{被动符号约定}的结果——两端口都用“流入网络”的速度记账（端口 1 用 $V_m$、端口 2 用 $-V_s$），从而端口功率 $P=\mathrm{Re}\{F V^*\}/2$ 流入网络为正，与本书能量端口 $P=\mathbf f^\top\mathbf v$（\cref{eq:os-port-power}）同号。
操作者阻抗记 $Z_h$、环境阻抗记 $Z_e$；二者均设为\emph{被动}（$\mathrm{Re}\{Z_h\},\mathrm{Re}\{Z_e\}\ge0$）。
\textbf{源约定换算}：原始遥操作文献（含本章源稿）多写 $[F_h;-V_s]=\mathbf H[V_m;F_e]$，与本书一致，无需换算；个别文献用缩放 $x_s=\lambda_p x_m$、$f_v=(u^2-v^2)/2$（波变量）等，将在 \cref{ch:teleop-wave} 出现时再注。
\end{note}

% ════════════════════════════════════════════════════════════════
\section{为何要“绝对”稳定}
\label{sec:ts-why}

\paragraph{动机：终端不是设计者能挑的。}
普通反馈控制里，被控对象的参数虽有不确定，但范围有限。双边遥操作不同：闭环的两端各接着一个\emph{不受控的活物}——人手与外部环境，二者的阻抗在巨大的区间里随时跳变。\cref{tab:ts-scenarios} 列出几种极端而真实的工况。

\begin{table}[H]\centering\small
\caption{遥操作终端阻抗的极端工况（master 端操作者阻抗 $Z_h$、slave 端环境阻抗 $Z_e$）}
\label{tab:ts-scenarios}
\begin{tabular}{lllp{4.4cm}}
\toprule
工况 & 操作者 $Z_h$ & 环境 $Z_e$ & 稳定性挑战 \\
\midrule
紧握 master + 撞刚墙 & 高（手臂刚度 $\sim 2000\,$N/m） & 极高（$>10^6\,$N/m） & 高环路增益 $\to$ 接触瞬间振荡 \\
松握 master + 自由空间 & 低（$\sim 50\,$N/m） & $0$ & 低阻尼 $\to$ 超调、极限环 \\
轻触 + 弹性物体 & 中 & 中 & 阻抗失配 $\to$ 能量缓慢积累 \\
突然放手 + 碰撞 & $0\to\infty$ & $0\to\infty$ & 瞬态切换 $\to$ 失稳 \\
\bottomrule
\end{tabular}
\end{table}

\paragraph{反面：名义稳定不够用。}
若只针对某一组“典型”$Z_h,Z_e$ 设计并验稳，系统在该工作点附近确实稳定——但操作者只要换个握法，或工具从自由空间一头撞上钢板，工作点就跳到表中另一行，名义稳定的保证当场失效。我们需要的是\emph{一次性覆盖整张表（及其任意中间组合）}的保证。

\begin{definition}[绝对稳定]\label{def:ts-abs-stable}
称由混合参数 $\mathbf H$ 描述的二端口网络是\emph{绝对稳定}的，若对\emph{任意}被动终端 $Z_h,Z_e$（$\mathrm{Re}\{Z_h\},\mathrm{Re}\{Z_e\}\ge0$），闭环系统都 BIBO 稳定。
\end{definition}

\paragraph{为何“被动终端”是恰当假设。}
把 $Z_h,Z_e$ 限制为被动，并非保守，而是物理事实的最弱合理刻画：人手与外部环境从端口看都不是能量源——它们可以储能、耗能，但不会凭空向系统注入无穷能量。\footnote{严格地说，肌肉可以做正功，但在感兴趣的频段与短时尺度上，把操作者建模为被动阻抗是遥操作稳定性分析的标准前提；主动操作者带来的“人在环”不稳定属另一专题。}于是“对任意被动终端稳定”恰好覆盖了 \cref{tab:ts-scenarios} 的整个条件空间。

\begin{insight}[Goertz 机械主从臂：免费的绝对稳定]
1950 年代 Goertz 在阿贡实验室造的机械主从臂——主从两端用钢丝绳和滑轮直接相连——同时拥有\emph{完美透明}与\emph{绝对稳定}。原因只有一个：纯被动机械系统没有电机、控制器、通信，\emph{根本没有创生能量的途径}。此后所有电控方案（二端口、波变量、TDPA）本质上都是\textbf{用主动控制器去模仿一个被动机械系统}。绝对稳定，就是给这种“模仿”立的及格线。
\end{insight}

% ════════════════════════════════════════════════════════════════
\section{Llewellyn 四判据}
\label{sec:ts-criteria}

绝对稳定的充要条件由 Llewellyn 在微波二端口理论中给出（1952），Lawrence（1993）将其引入遥操作\cite{lawrence1993stability}。它把“对一族无穷多终端都稳定”这一难以直接验证的要求，化为只关于 $\mathbf H$ 本身、可逐频检验的四个条件。

\begin{theorem}[Llewellyn 绝对稳定判据]\label{thm:ts-llewellyn}
设二端口网络由混合参数 $\mathbf H(s)$（\cref{eq:ts-hparam}）描述。该网络对任意被动终端绝对稳定（\cref{def:ts-abs-stable}），\emph{当且仅当}下列四条同时成立：
\begin{align}
  \text{(L1)}\quad & h_{11}(s),\,h_{22}(s)\ \text{在右半平面无极点（各自内部稳定）};\label{eq:ts-L1}\\
  \text{(L2)}\quad & \mathrm{Re}\{h_{11}(j\omega)\}\ge0,\quad\forall\omega;\label{eq:ts-L2}\\
  \text{(L3)}\quad & \mathrm{Re}\{h_{22}(j\omega)\}\ge0,\quad\forall\omega;\label{eq:ts-L3}\\
  \text{(L4)}\quad & \eta(\omega):=\frac{2\,\mathrm{Re}\{h_{11}\}\,\mathrm{Re}\{h_{22}\}-\mathrm{Re}\{h_{12}h_{21}\}}{|h_{12}h_{21}|}-1\ \ge0,\quad\forall\omega.\label{eq:ts-L4}
\end{align}
$\eta(\omega)$ 称\emph{稳定因子}（stability margin）。
\end{theorem}

\noindent 四条的物理读法如下（证明见 \cref{sec:ts-proof}）：

\begin{table}[H]\centering\small
\caption{Llewellyn 四判据的物理含义}
\label{tab:ts-criteria-meaning}
\begin{tabular}{cll}
\toprule
条件 & 数学要求 & 物理含义 \\
\midrule
L1 & $h_{11},h_{22}$ 各自稳定 & master 与 slave 子系统\emph{自身}内部稳定 \\
L2 & $\mathrm{Re}\{h_{11}\}\ge0$ & master 端口对操作者\emph{被动}（不向操作者注入能量） \\
L3 & $\mathrm{Re}\{h_{22}\}\ge0$ & slave 端口对环境\emph{被动}（不向环境注入能量） \\
L4 & $\eta(\omega)\ge0$ & 两端口耗散 $\ge$ 互耦合传递的能量（耦合不创生能量） \\
\bottomrule
\end{tabular}
\end{table}

\begin{insight}[为什么 L4 是灵魂]
L2、L3 只保证\emph{每个端口单独}被动。但两个各自被动的端口经反馈\emph{耦合}后，仍可能整体产生能量——正如一只运放（非被动）配上纯被动反馈网络就能起振：反馈网本身不产能，正反馈组合却能从电源汲能维持振荡。L4 正是给互耦合 $h_{12}h_{21}$ 设的上限，确保耦合不构成这种“正反馈起振”。L1--L3 都是 L4 之外的“边界必要条件”，L4 才是把它们补成\emph{充要}的那一块。
\end{insight}

\begin{note}
\textbf{$\eta$ 的几个等价/退化读数.}\;
当某频率 $h_{12}h_{21}=0$（端口完全解耦）时 \cref{eq:ts-L4} 的分母为零，此时 L4 自动满足（无互耦合即无起振之虞），实践中令 $\eta\to+\infty$ 处理。
$\eta\ge0$ 与不等式 $2\,\mathrm{Re}\{h_{11}\}\mathrm{Re}\{h_{22}\}-\mathrm{Re}\{h_{12}h_{21}\}\ge|h_{12}h_{21}|$ 完全等价；后者在数值上更稳健（避免除零），\cref{sec:ts-example} 与代码里都用它。
\end{note}

% ════════════════════════════════════════════════════════════════
\section{完整证明：从闭环特征方程到 \texorpdfstring{$\eta$}{eta}}
\label{sec:ts-proof}

理解 Llewellyn 判据不能止于记住四条——要看清它们\emph{如何}从“对任意被动终端稳定”一步步逼出来，才能在非标准架构里灵活套用。

\paragraph{思路。}
绝对稳定要求：对任意被动 $Z_h,Z_e$，闭环 BIBO 稳定，即闭环特征方程在右半平面无零点。我们先写出特征方程，再分别用“两端极端终端”逼出 L1--L3，用“一般被动终端”的端口功率非负逼出 L4。

\begin{derivation}[闭环特征方程]
操作者端口：操作者“推”进网络，其阻抗与被动符号相反，故
\begin{equation}
  F_h=-Z_h\,V_m.
  \label{eq:ts-master-term}
\end{equation}
环境端口：环境对 slave 呈阻抗，
\begin{equation}
  F_e=Z_e\,V_s.
  \label{eq:ts-env-term}
\end{equation}
把 \cref{eq:ts-master-term,eq:ts-env-term} 代入混合参数 \cref{eq:ts-hparam} 的两行：
\begin{align}
  -Z_h V_m &= h_{11}V_m+h_{12}Z_e V_s,\label{eq:ts-row1}\\
  -V_s &= h_{21}V_m+h_{22}Z_e V_s.\label{eq:ts-row2}
\end{align}
由 \cref{eq:ts-row2} 解出 $V_s(1+h_{22}Z_e)=-h_{21}V_m$，即 $V_s=-h_{21}V_m/(1+h_{22}Z_e)$。代回 \cref{eq:ts-row1}：
\begin{equation}
  (Z_h+h_{11})V_m=-h_{12}Z_e\cdot\frac{-h_{21}V_m}{1+h_{22}Z_e}
  =\frac{h_{12}h_{21}Z_e}{1+h_{22}Z_e}V_m.
\end{equation}
非平凡解（$V_m\neq0$）要求
\begin{equation}
  \boxed{\;(Z_h+h_{11})(1+h_{22}Z_e)=h_{12}h_{21}Z_e\;}
  \label{eq:ts-char}
\end{equation}
这就是闭环特征方程。BIBO 稳定 $\iff$ \cref{eq:ts-char} 对\emph{所有}被动 $Z_h,Z_e$ 在 $\mathrm{Re}(s)\ge0$ 上无根。
\end{derivation}

\paragraph{L1：两端开路/短路退化。}
取 $Z_h=0,Z_e=0$（自由空间 + 完全松手）。\cref{eq:ts-char} 退化为 $h_{11}\cdot1=0$，闭环极点即 $h_{11}$ 的极点；若 $h_{11}$ 有右半平面极点，响应发散。对 $h_{22}$ 同理（令 $Z_h\to\infty$ 配合分析）。故 \cref{eq:ts-L1} 必要。

\paragraph{L2、L3：一端开路、一端短路。}
取 $Z_e=0$（slave 端自由空间），\cref{eq:ts-char} 变为 $Z_h+h_{11}=0$，即 $Z_h=-h_{11}$。稳定要求对\emph{任意}被动 $Z_h$（$\mathrm{Re}\{Z_h\}\ge0$）都 $Z_h+h_{11}\neq0$ 于 $\mathrm{Re}(s)\ge0$。在 $s=j\omega$ 上，若 $\mathrm{Re}\{h_{11}(j\omega)\}<0$，则可取被动 $Z_h=-h_{11}(j\omega)$（其实部 $>0$，合法被动）使之恰好为零——出现纯虚根，临界失稳。故必须 $\mathrm{Re}\{h_{11}(j\omega)\}\ge0$，即 \cref{eq:ts-L2}。对称地令 $Z_h\to\infty$（操作者完全锁死）得 \cref{eq:ts-L3}。

\begin{derivation}[L4：一般被动终端 + 端口功率非负]
L1--L3 只锁定了终端的极端值。对一般 $Z_h=R_h+jX_h,\ Z_e=R_e+jX_e$（$R_h,R_e\ge0$），需更强条件。
在 $s=j\omega$ 处把 \cref{eq:ts-char} 的实部、虚部展开，要求对任意 $R_h,R_e\ge0$ 与任意 $X_h,X_e$，方程左右不能同时相等于一个纯虚根——等价地，收集关于 $R_hR_e,\ R_hX_e,\ X_hR_e,\ X_hX_e$ 的项并配方，可得稳定的充要不等式
\begin{equation}
  2\,\mathrm{Re}\{h_{11}\}\,\mathrm{Re}\{h_{22}\}-\mathrm{Re}\{h_{12}h_{21}\}\ \ge\ |h_{12}h_{21}|.
  \label{eq:ts-L4-raw}
\end{equation}
两边除以 $|h_{12}h_{21}|$（非零）再减 $1$，即得 $\eta(\omega)\ge0$（\cref{eq:ts-L4}）。\hfill$\blacksquare$
\end{derivation}

\paragraph{$\eta$ 的能量解读。}
上面的代数结论有一个干净的能量含义。在频率 $\omega$ 处，两端口注入网络的时均功率为
\begin{equation}
  P_1=\tfrac12\mathrm{Re}\{F_h V_m^*\},\qquad P_2=\tfrac12\mathrm{Re}\{F_e(-V_s)^*\}.
  \label{eq:ts-port-power}
\end{equation}
（这里的 $P=\tfrac12\mathrm{Re}\{F V^*\}$ 就是 \cref{eq:os-port-power} 的端口功率 $P=\mathbf f^\top\mathbf v$ 在正弦稳态、复幅值下的时均版本。）把 $\mathbf H$ 代入 \cref{eq:ts-port-power} 并要求\emph{网络耗散而非产能}，
\begin{equation}
  P_1+P_2\ge0\quad\text{对任意端口激励},
  \label{eq:ts-passive-sum}
\end{equation}
展开后正是 \cref{eq:ts-L4-raw}：分子里 $2\,\mathrm{Re}\{h_{11}\}\mathrm{Re}\{h_{22}\}$ 来自两端口自身的耗散项（$\mathrm{Re}\{h_{11}\}|V_m|^2$ 与 $\mathrm{Re}\{h_{22}\}|V_s|^2$），$\mathrm{Re}\{h_{12}h_{21}\}$ 与 $|h_{12}h_{21}|$ 来自互耦合传递。于是 $\eta\ge0$ $\iff$ 互联网络对任意激励无源——把 Llewellyn 判据接回了无源性这条本书主线（\cref{thm:os-kyp}）。

\begin{insight}[四判据各来自哪种终端]
\begin{itemize}
  \item L1 来自\emph{两端开路}（$Z_h=Z_e=0$）的退化；
  \item L2、L3 来自\emph{一端开路、一端短路}（$Z_e=0$ 或 $Z_h\to\infty$）；
  \item L4 来自\emph{两端都是一般被动阻抗}的完整情形，等价于端口功率 $P_1+P_2\ge0$。
\end{itemize}
四条互不可替代、合起来充要。若哪天有人声称找到“更弱的替代条件”，必是漏算了某类被动终端。
\end{insight}

% ════════════════════════════════════════════════════════════════
\section{Worked Example：1-DOF PD 位置跟踪 + 力反射}
\label{sec:ts-example}

把判据落成一次完整手算。考虑最常见的 PF 架构（位置 master$\to$slave，力 slave$\to$master）的 1 自由度模型。

\begin{table}[H]\centering\small
\caption{1-DOF worked example 的系统参数}
\label{tab:ts-example-param}
\begin{tabular}{lll}
\toprule
部件 & 参数 & 取值 \\
\midrule
Master 本体 & 质量 $M_m$、阻尼 $B_m$ & $0.5\,$kg，$2\,$Ns/m \\
Slave 本体 & 质量 $M_s$、阻尼 $B_s$ & $1.0\,$kg，$1\,$Ns/m \\
PD 位置跟踪 & $K_p,\ K_d$ & $100\,$N/m，$20\,$Ns/m \\
力反射增益 & $K_f$ & $0.8$ \\
Master 本地阻尼 & $B_{\rm loc}$ & $3\,$Ns/m \\
\bottomrule
\end{tabular}
\end{table}

\noindent slave 位置控制律 $\tau_s=K_p(x_m-x_s)+K_d(\dot x_m-\dot x_s)$；master 力反馈 $\tau_m=K_f f_e+B_{\rm loc}(-\dot x_m)$。这与机械臂关节动力学 $\mathbf M(\mathbf q)\ddot{\mathbf q}+\mathbf C(\mathbf q,\dot{\mathbf q})\dot{\mathbf q}+\mathbf g(\mathbf q)=\boldsymbol\tau$ 在单自由度、重力补偿后线性化的标量版一致；任务空间力 $f$ 与关节力矩 $\tau$ 经 $\tau=J^\top f$ 对偶（此处 $J=1$）。

\begin{derivation}[推导 $h$ 参数]
记 master 本体阻抗 $Z_m(s)=M_m s+B_m$、slave 本体阻抗 $Z_s(s)=M_s s+B_s$、PD 控制器 $C(s)=K_p+K_d s$。
\begin{itemize}
  \item $h_{11}=F_h/V_m\big|_{F_e=0}$：力反馈为零时 master 端只剩本体 + 本地阻尼，
  \[
    h_{11}(s)=Z_m(s)+B_{\rm loc}=M_m s+(B_m+B_{\rm loc})=0.5s+5.
  \]
  \item $h_{12}=F_h/F_e\big|_{V_m=0}$：master 锁死时环境力经反射增益直达 master 端力，
  \[
    h_{12}(s)=K_f=0.8.
  \]
  \item $h_{21}=-V_s/V_m\big|_{F_e=0}$：slave 位置闭环的速度跟踪传递函数，
  \[
    h_{21}(s)=-\frac{C(s)}{Z_s(s)+C(s)}=-\frac{20s+100}{21s+101}.
  \]
  \item $h_{22}=-V_s/F_e\big|_{V_m=0}$：master 锁死时 slave 闭环对外力的导纳，
  \[
    h_{22}(s)=\frac{1}{Z_s(s)+C(s)}=\frac{1}{21s+101}.
  \]
\end{itemize}
\end{derivation}

\paragraph{检查 L1--L3。}
$h_{11}=0.5s+5$ 无极点；$h_{22}=1/(21s+101)$ 极点在 $s=-101/21\approx-4.81$（左半平面）——L1 \checkmark。
$\mathrm{Re}\{h_{11}(j\omega)\}=5>0$——L2 \checkmark。
$\mathrm{Re}\{h_{22}(j\omega)\}=\dfrac{101}{101^2+(21\omega)^2}>0$——L3 \checkmark。

\begin{derivation}[计算 $\eta(\omega)$]
互耦合
\[
  h_{12}h_{21}=0.8\cdot\Bigl(-\frac{20s+100}{21s+101}\Bigr)=-\frac{16s+80}{21s+101}.
\]
在 $s=j\omega$，有理化分母：
\[
  \mathrm{Re}\{h_{12}h_{21}\}=-\frac{80\cdot101+16\cdot21\,\omega^2}{101^2+(21\omega)^2}
  =-\frac{8080+336\,\omega^2}{10201+441\,\omega^2}<0,
\]
（实部恒负，对 $\eta$ 有利），而
\[
  |h_{12}h_{21}|=\frac{0.8\sqrt{(20\omega)^2+100^2}}{\sqrt{(21\omega)^2+101^2}}.
\]
取直流 $\omega=0$：$\mathrm{Re}\{h_{11}\}=5$，$\mathrm{Re}\{h_{22}\}=101/10201\approx0.0099$，$\mathrm{Re}\{h_{12}h_{21}\}=-8080/10201\approx-0.792$，$|h_{12}h_{21}|=0.8\times100/101\approx0.792$。代入 \cref{eq:ts-L4}：
\[
  \eta(0)=\frac{2\times5\times0.0099-(-0.792)}{0.792}-1
  =\frac{0.099+0.792}{0.792}-1\approx0.125>0.
\]
取 $\omega=10\,$rad/s 同法可得 $\eta(10)\approx0.024>0$（裕度随 $\omega$ 升高而变薄，因 $\mathrm{Re}\{h_{22}\}\propto1/\omega^2$ 迅速衰减）。
\end{derivation}

\begin{insight}[结论与敏感性]
该系统在\emph{无时延}下满足 Llewellyn 绝对稳定：本地阻尼 $B_{\rm loc}=3\,$Ns/m 提供了足够耗散，使 $\eta>0$ 有正裕度。若去掉本地阻尼（$B_{\rm loc}=0$），$h_{11}=0.5s+2$、$\mathrm{Re}\{h_{11}\}=2$，$\eta$ 减小但仍可能为正——只是裕度变薄、对时延极其敏感（\cref{sec:ts-delay} 将看到时延如何吃光这点裕度）。\textbf{本地阻尼是买稳定性的“硬通货”，代价是透明度（\cref{sec:ts-tradeoff}）。}
\end{insight}

\noindent 把 $\eta(\omega)$ 的逐频检验写成代码只是几行——核心是\emph{用不等式形式 \cref{eq:ts-L4-raw} 避免除零}：

\begin{codebox}[Python]{逐频计算 Llewellyn 稳定因子 $\eta(\omega)$}
import numpy as np

def eta_of_omega(H, omega):           # H(s) -> 2x2 复矩阵的可调用对象
    s = 1j * omega
    h = H(s)
    h11, h12, h21, h22 = h[0,0], h[0,1], h[1,0], h[1,1]
    num = 2*h11.real*h22.real - (h12*h21).real
    den = abs(h12*h21)
    if den < 1e-12:                   # 端口解耦:L4 自动满足
        return np.inf
    return num/den - 1.0

def is_absolutely_stable(H, omegas):  # L2,L3,L4 逐频;L1 另查极点
    return all(H(1j*w)[0,0].real >= 0 and H(1j*w)[1,1].real >= 0
               and eta_of_omega(H, w) >= 0 for w in omegas)
\end{codebox}

% ════════════════════════════════════════════════════════════════
\section{透明度--稳定性的根本矛盾}
\label{sec:ts-tradeoff}

本章最重要、也最“扫兴”的结论：理想透明与绝对稳定在数学上互斥。

\begin{proposition}[理想透明度 $\Rightarrow$ 零鲁棒裕度]\label{prop:ts-ideal-zero}
对理想透明的混合参数 $\mathbf H_{\rm ideal}=\big[\begin{smallmatrix}0&1\\-1&0\end{smallmatrix}\big]$（$h_{11}=h_{22}=0,\ h_{12}=1,\ h_{21}=-1$，见 \cref{ch:teleop-twoport}），Llewellyn 判据恰好成立于边界，稳定因子
\begin{equation}
  \eta\equiv0.
  \label{eq:ts-eta-zero}
\end{equation}
\end{proposition}
\begin{proof}
逐条代入 \cref{thm:ts-llewellyn}：L1 中 $h_{11}=h_{22}=0$ 无极点；L2、L3 中 $\mathrm{Re}\{0\}=0\ge0$（恰在边界）。L4：
\[
  \eta=\frac{2\cdot0\cdot0-\mathrm{Re}\{1\cdot(-1)\}}{|1\cdot(-1)|}-1=\frac{0-(-1)}{1}-1=1-1=0.
\]
\end{proof}

\begin{insight}[$\eta=0$ 意味着什么]
$\eta=0$ 是“恰好及格”——名义上满足绝对稳定，但\emph{任何}把 $\eta$ 推向负的微小扰动都会失稳。而这些扰动在真实系统里无可避免：
\begin{table}[H]\centering\small
\caption{把理想透明系统推下稳定边界的现实扰动}
\label{tab:ts-perturb}
\begin{tabular}{ll}
\toprule
扰动来源 & 对 $\eta$ 的影响 \\
\midrule
通信时延 $T>0$ & 引入 $e^{-sT}$，改变 $\mathrm{Re}\{h_{12}h_{21}\}$，$\eta$ 可转负（\cref{sec:ts-delay}） \\
电机惯量/相位 & $h_{11}\neq0$ 且带相位，某些频率 $\mathrm{Re}\{h_{11}\}<0$ \\
传感器噪声 & 等效注入能量，破坏端口被动 \\
量化/离散 & 引入非线性，LTI 的 Llewellyn 分析失效 \\
\bottomrule
\end{tabular}
\end{table}
\end{insight}

\paragraph{解药与代价：加本地阻尼。}
要换得正的裕度，物理上就是在 master、slave 端各加一点本地阻尼，等价于把对角元抬离零：
\begin{equation}
  h_{11}\to h_{11}+b_m,\qquad h_{22}\to h_{22}+b_s,\qquad b_m,b_s>0.
  \label{eq:ts-add-damp}
\end{equation}
保持理想的反对角 $h_{12}h_{21}=-1$ 不变，代入 \cref{eq:ts-L4}：

\begin{derivation}[加阻尼后的稳定因子 $\eta'$]
\begin{equation}
  \eta'=\frac{2\,b_m\,b_s-\mathrm{Re}\{-1\}}{|-1|}-1=\frac{2b_mb_s+1}{1}-1=2\,b_m\,b_s\ >\ 0.
  \label{eq:ts-eta-prime}
\end{equation}
\end{derivation}

\noindent\textbf{只要 $b_m,b_s>0$（无论多小），$\eta'>0$。} 但天下没有免费的稳定：操作者感受阻抗（\cref{ch:teleop-twoport} 的 $Z_{to}$）不再等于真实环境 $Z_e$——多出来的对角阻尼让操作者凭空感到一层“黏滞”，即 $Z_{to}\neq Z_e$。透明度被牺牲了。

\begin{insight}[两种视角看这条 trade-off]
\begin{itemize}
  \item \textbf{控制论视角}：这是 Bode 灵敏度积分约束的力学版——“阻抗匹配在某处改善，必在鲁棒裕度处付费”，就像灵敏度函数在某频段压低必在别处抬高。
  \item \textbf{热力学视角}：无阻尼的理想遥操作系统是“可逆”的（$\eta=0$），如卡诺机效率最高却不可实现；加阻尼把它变“不可逆”（$\eta>0$），牺牲效率（透明度）换可靠（稳定）。
\end{itemize}
\end{insight}

\begin{pitfall}[思维陷阱：“提高控制带宽就能提高透明度”]
直觉以为“更快的控制器 $\to$ 更好的跟踪 $\to$ 更透明”。但：(a) 高带宽更易撞上增益裕度而失稳；(b) 采样-保持在高频引入额外相位滞后；(c) 传感器噪声在高频被放大。透明度-稳定性 trade-off \emph{不是}“换个更好的控制器”能消除的——它是系统级的物理限制。同理，单纯增大力反馈增益 $|h_{12}|$ 会增大 $|h_{12}h_{21}|$，使 \cref{eq:ts-L4} 更难满足，透明不升反降。
\end{pitfall}

% ════════════════════════════════════════════════════════════════
\section{时延的定量影响}
\label{sec:ts-delay}

前面所有分析都在“无时延”下进行。一旦 master 与 slave 之间有通信时延 $T$，局面急转直下。

\paragraph{时延把延迟因子塞进互耦合。}
设无时延时位置/力两通道分别为 $h_{21}(s)=-C_p(s)$（位置跟踪）、$h_{12}(s)=C_f(s)$（力反射）。加入往返时延 $T$（每方向 $T/2$），两通道各乘一个延迟因子：
\begin{equation}
  h_{12}(s)=C_f(s)\,e^{-sT/2},\qquad h_{21}(s)=-C_p(s)\,e^{-sT/2}.
  \label{eq:ts-delay-channels}
\end{equation}
于是在 $s=j\omega$ 处，\emph{互耦合乘积}恰好吃满整个往返延迟：
\begin{equation}
  h_{12}(j\omega)h_{21}(j\omega)=-|C_fC_p|\,e^{-j\omega T},\qquad
  \mathrm{Re}\{h_{12}h_{21}\}=-|C_fC_p|\cos(\omega T).
  \label{eq:ts-delay-recoupling}
\end{equation}

\begin{derivation}[临界频率 $\omega=\pi/T$：余弦翻号]
注意 \cref{eq:ts-delay-recoupling} 里 $\mathrm{Re}\{h_{12}h_{21}\}$ 随 $\cos(\omega T)$ 振荡。在 $\omega T=\pi$（即 $\omega=\pi/T$）处，$\cos\pi=-1$，本来“有利”的负实部翻成\emph{正}：
\[
  \mathrm{Re}\{h_{12}h_{21}\}\big|_{\omega=\pi/T}=-|C_fC_p|\cos\pi=+|C_fC_p|.
\]
对加了本地阻尼的对角元 $\mathrm{Re}\{h_{11}\}\approx b_m,\ \mathrm{Re}\{h_{22}\}\approx b_s$，代入 \cref{eq:ts-L4}：
\begin{equation}
  \eta\Bigl(\tfrac{\pi}{T}\Bigr)=\frac{2b_mb_s-|C_fC_p|}{|C_fC_p|}-1=\frac{2b_mb_s}{|C_fC_p|}-2.
  \label{eq:ts-eta-critical}
\end{equation}
要 $\eta(\pi/T)\ge0$，须
\begin{equation}
  \boxed{\;b_m\,b_s\ \ge\ |C_fC_p|\;}
  \label{eq:ts-damp-requirement}
\end{equation}
即\emph{两端本地阻尼之积，必须压过两通道增益之积}。
\end{derivation}

\paragraph{坏消息：时延不改阻尼量级，却把临界频率拖进操作频段。}
\cref{eq:ts-damp-requirement} 与 $T$ 无关——对固定增益，所需阻尼量级不随时延变。真正致命的是\emph{临界频率} $\omega=\pi/T$ 随 $T$ 增大而\emph{降低}：当它落进人手操作频段（约 $1$--$10\,$Hz，即 $\sim6$--$63\,$rad/s），系统在\emph{正常操作}中就会失稳。\cref{tab:ts-delay-table} 给出 $|C_fC_p|=1$ 时的临界频率。

\begin{table}[H]\centering\small
\caption{时延 $T$ 与 Llewellyn 临界频率 $\omega=\pi/T$（设 $|C_fC_p|=1$）}
\label{tab:ts-delay-table}
\begin{tabular}{lll}
\toprule
时延 $T$ & 临界频率 $\pi/T$ & 后果 \\
\midrule
$1\,$ms & $3142\,$rad/s & 远超操作频段，无影响 \\
$10\,$ms & $314\,$rad/s & 仍在操作频段之上 \\
$100\,$ms & $31.4\,$rad/s（$\approx5\,$Hz） & 已落入手操作频段上沿，临界 \\
$200\,$ms & $15.7\,$rad/s（$\approx2.5\,$Hz） & 正常操作中即失稳 \\
$500\,$ms & $6.3\,$rad/s（$\approx1\,$Hz） & 手操作主频段失稳 \\
$1000\,$ms & $3.14\,$rad/s（$\approx0.5\,$Hz） & 连缓慢操作都不稳 \\
\bottomrule
\end{tabular}
\end{table}

\begin{insight}[$>200\,$ms：Llewellyn 几近无解，必须换框架]
当 $T\gtrsim200\,$ms，临界频率压进人手主操作频段，要靠 \cref{eq:ts-damp-requirement} 硬撑就得加到\emph{操作者根本推不动}的阻尼——透明度归零，系统名存实亡。深空遥操作（地月往返 $\sim2.6\,$s；地火单向光行时 $3$--$22\,$min，往返加倍）因此\emph{彻底放弃} bilateral 力反馈，改走监督控制。要在大时延下\emph{既稳又有力反馈}，必须跳出“逐频压 $\eta$”的思路，改用\emph{编码能量传输}的波变量/散射方法——它对\emph{任意常数时延}都保证通信通道无源。这正是下一章（\cref{ch:teleop-wave}）的全部动机。
\end{insight}

\begin{pitfall}[概念误区：“Llewellyn 太保守，实际不需要”]
有人嫌 Llewellyn“要求任意被动负载下稳定，过严”。但它是\emph{充要}条件，不是仅充分：一旦违反任一条，\emph{必}存在某个被动 $Z_h,Z_e$ 使系统失稳。而人手刚度在 $50$--$5000\,$N/m（松握$\to$紧握）、环境从自由空间（$Z_e=0$）到刚墙（$Z_e\to\infty$）——真实工况\emph{完整覆盖}了“任意被动负载”的条件空间。所以这里没有“保守”的余地：判据失败处就是真实事故的所在。
\end{pitfall}

% ════════════════════════════════════════════════════════════════
\section{\texorpdfstring{$Z$-width}{Z-width} 与采样率：复用 Colgate--Schenkel 界}
\label{sec:ts-zwidth}

前几节的 $\eta$ 都在连续时间下推。但所有遥操作器都是数字采样实现的，离散化本身会破坏无源——这一关本书在操作空间一章已经证过，这里直接\emph{升级复用}，不重证。

\paragraph{离散化破坏无源。}
零阶保持给回路引入相位滞后 $e^{-j\omega T/2}$（$T$ 为采样周期），把一个名义无源的虚拟阻抗变成有源元件。Colgate 与 Schenkel（1997）给出“虚拟墙”采样数据无源的充分条件\cite{colgate1997passivity}：物理阻尼 $B$、虚拟阻尼 $D_d$ 与虚拟刚度须满足本书已证的
\begin{equation}
  K_d<\frac{2(B-D_d)}{T}\quad\Longleftrightarrow\quad
  K_{\max}=\frac{2(B-D_d)}{T},
  \tag{\ref{eq:os-discrete-passive} 复用}
\end{equation}
即\cref{eq:os-discrete-passive}。结论一句话：\emph{采样越慢，可被动渲染的刚度上限越低}。（注：虚拟阻尼 $D_d$ 以\textbf{减号}进入——这是 Colgate--Schenkel 的反直觉点：ZOH 下用差分估计速度带一拍延迟，使虚拟阻尼在采样系统中反而\emph{注能}、须由物理阻尼 $B$ 额外补偿，故 $K_{\max}=2(B-D_d)/T$；$b>KT/2+B$ 为无源充要条件，$b>KT/2-B$ 则是更弱的\emph{稳定}条件。本章数值均取 $D_d\to0$、$K_{\max}=2B/T$，不受此符号影响。详见母式 \cref{eq:os-discrete-passive} 处推导。）

\begin{definition}[$Z$-width]\label{def:ts-zwidth}
触觉/遥操作系统的 $Z$-width 是其在保持被动（不向操作者注入能量）前提下能渲染的阻抗范围 $[Z_{\min},Z_{\max}]$：$Z_{\min}$ 越低（自由空间越“空”）、$Z_{\max}$ 越高（硬墙越“硬”）越好。由 \cref{eq:os-discrete-passive}，可渲染的最大刚度 $K_{\max}=2(B-D_d)/T$ 直接由采样周期 $T$ 与可用阻尼决定。
\end{definition}

\begin{insight}[力反馈可行性的“第一道门槛”]
设计力反馈遥操作时，$Z$-width 是\emph{是否值得做}的前置判断，可走一张决策图：
\begin{enumerate}
  \item 测物理阻尼 $b$ 与采样周期 $T$；
  \item 算 $K_{\max}=2b/T$（取 $D_d\to0$ 的保守估计）；
  \item 比对目标环境刚度 $K_{\rm env}$；
  \item $K_{\max}\ge K_{\rm env}$？\;是 $\Rightarrow$ 可做力反馈，继续选架构；否 $\Rightarrow$ 要么提采样率（硬件升级），要么放弃力反馈。
\end{enumerate}
\end{insight}

\paragraph{为何 ALOHA 在 $50\,$Hz 下不做力反馈。}
以 ALOHA 这类低成本遥操作（约 $50\,$Hz 控制频率、Python 软实时、无力传感器）为例：$T=1/50=0.02\,$s，取物理阻尼 $b\approx2\,$Ns/m，则
\begin{equation}
  K_{\max}=\frac{2b}{T}=\frac{2\times2}{0.02}=200\ \text{N/m}.
  \label{eq:ts-aloha-zwidth}
\end{equation}
而人手抓握刚度 $\sim2000\,$N/m——$K_{\max}=200\,$N/m \emph{连一面像样的软墙都渲染不出}，更别说硬接触。结论：在 $50\,$Hz 采样下，被动力反馈\emph{物理上不可行}，瓶颈是控制频率与采样-保持效应，而非传感器成本。ALOHA 因此理性地退化为\emph{单向位置映射}（只开 master$\to$slave 位置通道），把省下的带宽与可靠性留给模仿学习的数据采集——这是系统级 co-design 的结果，不是“省钱”。

\begin{insight}[采样率决定力反馈质量]
\cref{eq:os-discrete-passive} 揭示一条贯穿触觉与遥操作的硬约束：\emph{力反馈的“硬度”上限正比于采样率}。工程经验是采样率取力控带宽的 $10$ 倍以上；要渲染 $2000\,$N/m 的刚墙、保持 $b\approx2\,$Ns/m，需 $T\le2\times2/2000=2\,$ms，即 $\ge500\,$Hz——这正是专业力反馈触觉设备（如 Phantom、Force Dimension omega/sigma 系列）的伺服回路跑在 $1$--$4\,$kHz 量级的根本原因。\pz{反例校正：原稿举“da Vinci 操控台”为力反馈触觉设备实例有误——经典 da Vinci（Si/Xi）以\emph{缺少}力/触觉反馈著称，仅 2024 年的 da Vinci 5 才加入有限的器械尖端测力（且尚非连续力反射）；故此处改用 Phantom 与 Force Dimension omega/sigma 这类真正的接地式动觉力反馈设备。Omega.7 刷新率实测可达 $4\,$kHz，因此写作 $1$--$4\,$kHz 量级而非单一 $1\,$kHz。}
\end{insight}

% ════════════════════════════════════════════════════════════════
\section{本章小结}
\label{sec:ts-summary}

\begin{itemize}
  \item \textbf{绝对稳定}（\cref{def:ts-abs-stable}）是双边遥操作的及格线：对任意被动操作者 $Z_h$ 与环境 $Z_e$ 都不失稳，因为终端是不受控的活物，会在 \cref{tab:ts-scenarios} 的极端工况间任意游走。
  \item \textbf{Llewellyn 四判据}（\cref{thm:ts-llewellyn}）把这一要求化为只关于 $\mathbf H$ 的可逐频检验：L1（各自稳定）、L2/L3（两端口被动）、L4（$\eta(\omega)\ge0$，互耦合不创生能量）。其中 L4 是灵魂，等价于端口功率 $P_1+P_2\ge0$（\cref{eq:ts-passive-sum}），接回无源性主线。
  \item \textbf{证明}的钥匙是闭环特征方程 $(Z_h+h_{11})(1+h_{22}Z_e)=h_{12}h_{21}Z_e$（\cref{eq:ts-char}）：极端终端逼出 L1--L3，一般被动终端逼出 L4。
  \item \textbf{透明度-稳定性根本矛盾}：理想透明 $\Rightarrow\eta\equiv0$（\cref{prop:ts-ideal-zero}，零鲁棒裕度）；加本地阻尼 $\Rightarrow\eta'=2b_mb_s>0$（\cref{eq:ts-eta-prime}）但 $Z_{to}\neq Z_e$——稳定要花透明度买。
  \item \textbf{时延}把 $e^{-j\omega T}$ 塞进互耦合（\cref{eq:ts-delay-recoupling}）：在 $\omega=\pi/T$ 处余弦翻号、$\eta$ 转负，阻尼要求升为 $b_mb_s\ge|C_fC_p|$（\cref{eq:ts-damp-requirement}）；时延把临界频率拖进人手操作频段，$>200\,$ms 时 Llewellyn 几近无解——引出波变量（\cref{ch:teleop-wave}）。
  \item \textbf{$Z$-width 与采样率}：离散化破坏无源，可渲染刚度上限 $K_{\max}=2(B-D_d)/T$ 复用本书 Colgate--Schenkel 界（\cref{eq:os-discrete-passive}）；ALOHA $50\,$Hz $\Rightarrow K_{\max}\approx200\,$N/m，故力反馈不可行。
\end{itemize}

\begin{table}[H]\centering\small
\caption{本章符号表}
\label{tab:ts-symbols}
\begin{tabular}{lll}
\toprule
符号 & 含义 & 出处 \\
\midrule
$\mathbf H,\ h_{ij}$ & 混合参数矩阵及其元 & \cref{eq:ts-hparam} \\
$F_h,V_m$ & master（操作者侧）端口力/速度 & \cref{eq:ts-hparam} \\
$F_e,V_s$ & slave（环境侧）端口力/速度 & \cref{eq:ts-hparam} \\
$Z_h,Z_e$ & 操作者/环境阻抗（被动） & \cref{def:ts-abs-stable} \\
$Z_{to}$ & 操作者感受阻抗（透明度量度） & \cref{ch:teleop-twoport} \\
$\eta(\omega)$ & Llewellyn 稳定因子 & \cref{eq:ts-L4} \\
$C_p,C_f$ & 位置跟踪/力反射通道增益 & \cref{eq:ts-delay-channels} \\
$b_m,b_s$ & master/slave 本地阻尼 & \cref{eq:ts-add-damp} \\
$T$ & 通信往返时延 / 采样周期 & \cref{eq:ts-delay-recoupling}，\cref{eq:os-discrete-passive} \\
$K_{\max}$ & 被动可渲染最大刚度（$Z$-width） & \cref{def:ts-zwidth} \\
\bottomrule
\end{tabular}
\end{table}

\begin{practice}[本章练习]
\begin{enumerate}
  \item \textbf{[PP 架构判稳]} 对 1-DOF PP 架构（虚拟弹簧 $K_m$ 耦合 master--slave，两端各加本地阻尼 $b_m,b_s$），推导 $h_{11},h_{12},h_{21},h_{22}$，写出 $\eta(\omega)$，并求使 $\eta(\omega)\ge0\,\forall\omega$ 的最大 $K_m$。
  \item \textbf{[复算 worked example]} 用 \cref{tab:ts-example-param} 参数，把 $\eta(\omega)$ 在 $\omega\in[0,100]\,$rad/s 上画出；再令 $B_{\rm loc}=0$ 重画，比较裕度变化，并指出对时延的敏感性。
  \item \textbf{[时延极限]} 对深空遥操作 $T=2.6\,$s（地月往返），用 \cref{eq:ts-damp-requirement} 与 \cref{tab:ts-delay-table} 分析：临界频率落在何处？要靠加阻尼维持 Llewellyn 稳定，需多大的 $b_mb_s$？由此论证为何必须放弃 bilateral 力反馈。
  \item \textbf{[$Z$-width 设计]} 要在 $b=2\,$Ns/m 下被动渲染 $K_{\rm env}=3000\,$N/m 的刚墙，用 \cref{eq:os-discrete-passive} 求所需最小采样率；若硬件只能跑 $200\,$Hz，可渲染的最硬墙是多少？
  \item \textbf{[思考题]} 为什么 Lawrence 用阻抗 $Z_{to}$（而非力跟踪误差 $|f_h-f_e|$）定义透明度？提示：考虑缩放系统 $f_h=10f_e,\ v_s=v_m/10$——力误差很大，但操作者的\emph{阻抗感受}如何？
\end{enumerate}
\end{practice}
